\section{Introduction}\label{sec2}

Let $\mathcal{A} $ be an abelian category. Homological algebra is generally considered to be the study of chain complexes up to \emph{homology}, hence the name. This means that if two chain complexes $C_{\bullet} ,D_{\bullet} $ have the same homology $H_i(C_{\bullet} )\cong H_i(D_{\bullet} )$ for all $i$, then we should consider $C_{\bullet} $ and $D_{\bullet} $ to be the the same. Put another way, if the homology chain complexes $H_{\bullet} (C_{\bullet} )\cong H_{\bullet} (D_{\bullet} )$ are isomorphic, then we would like it to be true that $C_{\bullet} \cong D_{\bullet} $. Unfortunately, this is \emph{not true in general}.

This is a much bigger deal than the reader is likely aware of. In nearly all of math, we are interested in studying objects up to isomorphism. The objects of study in homological algebra are chain complexes, so we would think that we are studying them up to isomorphism. However, we \emph{are not} studying them up to isomorphism, but rather up to homology. This means that all of the methods we have developed to study objects up to isomorphism are now useless.

The easiest way to fix this issue is to define a \emph{new} category, called the \emph{derived category}. A \emph{quasi-isomorphism} is a morphism of chain complexes $f_{\bullet} : C_{\bullet}  \to D_{\bullet} $ which induces an isomorphism on their homology:
\[
	H_{\bullet} (f_{\bullet} ): H_{\bullet} (C_{\bullet} )\cong H_{\bullet} (D_{\bullet} ) \iff H_i(f_{\bullet} ) : H_i(C _{\bullet} )\cong H_i(D_{\bullet} )
.\]
This means that the study of chain complexes up to homology is the same as the study of chain complexes up to quasi-isomorphism. If two chain complexes $C_{\bullet} ,D_{\bullet} $ are quasi-isomorphic, they should be considered the same. In order to do this, we construct a new category $\mathcal{D} (\mathcal{A} )$ called the \emph{derived category} of $\mathcal{A} $, which has the same objects as $\Ch(\mathcal{A} ) $, but the quasi-isomorphisms have somehow been turned into isomorphisms. This turns the study of chain complexes in $\Ch(\mathcal{A} ) $ up to quasi-isomorphism into the study of chain complexes in $\mathcal{D} (\mathcal{A} )$ up to isomorphism.

As we will see in \cref{sec2}, the derived category oftentimes introduces more problems than it solves. Working in the derived category can be quite complex, meaning we have not considerably improved our situation. Modern homotopical algebra provides a more elegant solution: a \emph{derived $\infty $-category} $\mathcal{D} _{\infty } (\mathcal{A} )$.

An $(\infty ,1)$-category, which we will simply call an $\infty $-category, is like a normal category (which we call a 1-category) except it contains \emph{higher homotopical information}. An $\infty $-category $\mathcal{C} $ has a class $\mathcal{C} _0$ of objects, $\mathcal{C} _1$ of morphisms (sometimes called 1-morphisms), but composition of morphisms is not well-defined. To fix this, it also contains a class $\mathcal{C} _2$ of \emph{2-morphisms} or \emph{homotopies}, which are morphisms between the 1-morphisms. So if $f,g : x \to y$ are 1-morphisms, then we can talk about a homotopy $\alpha : f \to g$. The reason we call the homotopies is because they are required to be invertible. We can then define an equivalence relation on the set of 1-morphisms from $x$ to $y$, denoted $\mathcal{C} _1(x,y)$, by $f\sim g$ if there is a homotopy $f \to g$. Composition of 1-morphisms is well-defined up to homotopy in this sense.

We then have 3-morphisms, which are homotopies of homotopies, which must also be invertible. We then have 4-morphisms, which are homotopies of homotopies of homotopies, which must also be invertible, and so on to infinity. Explicitly describing these higher homotopies is very difficult, so usually if we need to go above 2-morphisms or 3-morphisms, we instead use a \emph{mapping object}. Given objects $x,y\in \mathcal{C} _0$, we define $\mathcal{C} (x,y)$ as a special type of $\infty $-category where \emph{all} levels of morphisms are invertible, called an $\infty $-groupoid. Objects are 1-morphisms, so $\mathcal{C} (x,y)_0=\mathcal{C} _1(x,y)$. Morphisms are homotopies, so $\mathcal{C} (x,y)_1=\mathcal{C} _2(x,y)$, and so on:
\[
	\mathcal{C} (x,y)_n=\mathcal{C} _{n+1} (x,y)
.\]
We will then define composition as some kind of morphism of these mapping objects, and will treat the higher homotopical information in this way.

This informal description of $\infty $-categories will suffice for our talk. A truly rigorous treatement with a precise definition of $\infty $-categories would require knowledge of simplicial sets, followed by a short lecture series worth of information about \emph{quasi-categories}. The interested reader should see \cite{HTT}. We will proceed with the informal description given here.
